\clearpage\begingroup
\HistoricalHead{184}{L}{PHYSIK.\quad NACHLASS.}
\begin{center}6.\end{center}
\noindent\hspace*{\parindent}Haben wir anstatt eines Rechtecks, eine grosse Anzahl gleicher und in gleichen Distanzen auf der Oberfläche eines rechtwinklichten Prisma liegender Vierecke, die alle von gleichen Strömen in übereinstimmender Richtung durchlaufen werden, so lässt sich die Summe aller Kräfte hinreichend genau durch Integration finden. Für alle diese Rechtecke haben nemlich $\alpha,\beta,\gamma,\delta$ gleiche Werthe, während $k$ ungleiche in arithmetischer Progression haben wird; ist $n$ die Anzahl der Rechtecke und $\lambda$ die Entfernung zweier auf einander folgender, und die äussersten Werthe von $k$ diese $k^0+\frac12\lambda$, $k'-\frac12\lambda$, so dass $k'=k^0+n\lambda$, so wird man die Summe aller Werthe jeder der drei Kräfte hinlänglich genau finden, wenn man ihren Ausdruck, $k$ als veränderlich betrachtend, mit $dk$ multiplicirt, von $k=k^0$ bis $k=k'$ integrirt und das gefundene Integral mit $\lambda$ dividirt.

Betrachten wir zuerst dieses Integral in Beziehung auf die beiden ersten Theile der ersten Kraft, $\frac{\alpha\delta}{aae}+\frac{\alpha\delta}{dde}$. Hier wird solches
\[
=\int\frac{\alpha\delta}{\surd(\alpha\alpha+\delta\delta+kk)}\left(\frac{1}{\alpha\alpha+kk}+\frac{1}{\delta\delta+kk}\right)dk.
\]
\noindent Dies Integral wird von $k=0$ anfangend
\[
\begin{aligned}
&=\operatorname{Arc\,sin}\frac{\delta k}{\surd(\alpha\alpha+\delta\delta)(\alpha\alpha+kk)}+\operatorname{Arc\,sin}\frac{\alpha k}{\surd(\alpha\alpha+\delta\delta)(\delta\delta+kk)}\\[5pt]
&=\operatorname{Arc\,tg}\frac{\delta k}{\alpha\surd(\alpha\alpha+\delta\delta+kk)}+\operatorname{Arc\,tg}\frac{\alpha k}{\delta\surd(\alpha\alpha+\delta\delta+kk)}
\end{aligned}
\]
\clearpage\SourceLeaf{190}\HistoricalHead{185}{R}{WECHSELWIRKUNGEN ZWISCHEN GALVANISCHEN STRÖMEN UND MAGNETISMUS.}
\[
\begin{aligned}
&=\operatorname{Arc\,tg}\frac{k\surd(\alpha\alpha+\delta\delta+kk)}{\alpha\delta}\\[5pt]
&=\operatorname{Arc\,cos}\frac{\alpha\delta}{\surd(kk+\alpha\alpha)(kk+\delta\delta)}.
\end{aligned}
\]
\noindent Also innerhalb der Grenzen der Integration
\[
=\operatorname{Arc\,tg}\frac{k'\surd(\alpha\alpha+\delta\delta+k'k')}{\alpha\delta}-\operatorname{Arc\,tg}\frac{k^0\surd(\alpha\alpha+\delta\delta+k^0k^0)}{\alpha\delta}.
\]
\noindent Wenn man nach der Analogie $\alpha\alpha+\delta\delta+k^0k^0=e^0e^0$ setzt u. s. w., so wird die ganze Kraft I
\[
\begin{aligned}
={}&\frac{1}{\lambda}\biggl\{\operatorname{Arc\,tg}\frac{k'e'}{\alpha\delta}+\operatorname{Arc\,tg}\frac{k'f'}{\alpha\beta}+\operatorname{Arc\,tg}\frac{k'g'}{\beta\gamma}+\operatorname{Arc\,tg}\frac{k'h'}{\gamma\delta}\\[5pt]
&\phantom{\frac{1}{\lambda}}-\operatorname{Arc\,tg}\frac{k^0e^0}{\alpha\delta}-\operatorname{Arc\,tg}\frac{k^0f^0}{\alpha\beta}-\operatorname{Arc\,tg}\frac{k^0g^0}{\beta\gamma}-\operatorname{Arc\,tg}\frac{k^0h^0}{\gamma\delta}\biggr\}.
\end{aligned}
\]
\noindent Ebenso entsteht aus einem Theile der Kraft II das Integral
\[
\int\frac{\delta k\mathbin{.}dk}{(\alpha\alpha+kk)\surd(\alpha\alpha+\delta\delta+kk)}=-\log\left(\frac{\delta+\surd(\alpha\alpha+\delta\delta+kk)}{\surd(\alpha\alpha+kk)}\right)+\text{Const.},
\]
\noindent also innerhalb der Grenzen der Integration
\[
\log\frac{\delta+\surd(\alpha\alpha+\delta\delta+k'k')}{\delta+\surd(\alpha\alpha+\delta\delta+k^0k^0)}\left[=\log\frac{\delta+e'}{\delta+e^0}\right].
\]
\noindent Ähnliche Ausdrücke geben die andern drei Theile und die vier Theile von III. Es wird folglich die ganze Kraft II
\[
=\frac{1}{\lambda}\log\left\{\frac{\delta+e'}{\delta+e^0}\mathbin{.}\frac{\beta+f'}{\beta+f^0}\mathbin{.}\frac{\beta+g^0}{\beta+g'}\mathbin{.}\frac{\delta+h^0}{\delta+h'}\right\}.
\]
\noindent Und auf gleiche Weise die ganze Kraft III
\[
=\frac{1}{\lambda}\log\left\{\frac{\alpha+e^0}{\alpha+e'}\mathbin{.}\frac{\alpha+f'}{\alpha+f^0}\mathbin{.}\frac{\gamma+g'}{\gamma+g^0}\mathbin{.}\frac{\gamma+h^0}{\gamma+h'}\right\}.
\]
\noindent Es ist überflüssig zu bemerken, dass in diesen Formeln die Logarithmen hyperbolische und dass die Bögen in Theilen des Halbmessers auszudrücken sind.
\clearpage\endgroup
